\SourcePage{371}{374}
\section{Relation between molecular and atomic terms}

As the distance between the nuclei of a diatomic molecule is increased, the limiting result is a pair of isolated atoms (or ions). This raises the question of how an electronic term of the molecule corresponds to the states of the atoms obtained when they are separated (E.~Wigner, E.~Witmer, 1928). The correspondence is not unique: bringing together two atoms in specified states may produce a molecule in any of several electronic states.

First suppose that the molecule contains two unlike atoms.
Let the isolated atoms occupy states with orbital angular momenta $L_1$, $L_2$ and spins $S_1$, $S_2$, where $L_1\geqslant L_2$.
Their projections on the line joining the nuclei take the values $M_1=-L_1,-L_1+1,\ldots,L_1$ and $M_2=-L_2,-L_2+1,\ldots,L_2$.
As the atoms approach, the resulting angular momentum $\Lambda$ is fixed by $\Lambda=|M_1+M_2|$.
Combining every possible pair of $M_1$ and $M_2$, we obtain the following frequencies for the distinct values of $\Lambda$:

\[
\begin{array}{c@{\qquad}l}
\Lambda=L_1+L_2&2\text{ times},\\
L_1+L_2-1&4\text{ times},\\
\dotfill&\dotfill\\
L_1-L_2&2(2L_2+1)\text{ times},\\
L_1-L_2-1&2(2L_2+1)\text{ times},\\
\dotfill&\dotfill\\
1&2(2L_2+1)\text{ times},\\
0&2L_2+1\text{ times}.
\end{array}
\]
Terms with $\Lambda\ne0$ are doubly degenerate, whereas those with $\Lambda=0$ are nondegenerate. It follows that the possible results are
\begin{equation}
\begin{array}{r@{\quad}c@{\quad}l}
1\text{ term}&\text{with}&\Lambda=L_1+L_2,\\
2\text{ terms}&\text{with}&\Lambda=L_1+L_2-1,\\
\dotfill&&\dotfill\\
2L_2+1\text{ terms}&\text{with}&\Lambda=L_1-L_2,\\
2L_2+1\text{ terms}&\text{with}&\Lambda=L_1-L_2-1,\\
\dotfill&&\dotfill\\
2L_2+1\text{ terms}&\text{with}&\Lambda=0;
\end{array}
\tag{80.1}
\end{equation}
altogether there are $(2L_2+1)(L_1+1)$ terms, whose values of $\Lambda$ range from 0 to $L_1+L_2$.

The spins $S_1$, $S_2$ of the two atoms combine, by the general rule for addition of angular momenta, into the total molecular spin and give the following

\SourcePage{372}{375}
possible values of $S$:
\begin{equation}
 S=S_1+S_2,\quad S_1+S_2-1,\ldots,|S_1-S_2|.
 \tag{80.2}
\end{equation}
Combining each of these values with all the values of $\Lambda$ in (80.1) gives the complete list of possible terms of the molecule formed.

For the $\Sigma$ terms, their sign remains to be determined. This can readily be done by observing that, as $r\to\infty$, the molecular wavefunctions can be written as products (or sums of products) of the wavefunctions of the two atoms. The value $\Lambda=0$ can arise either by adding nonzero projections $M_1=-M_2$, or when $M_1=M_2=0$. Denote the wavefunctions of the first and second atoms by $\psi_{M_1}^{(1)}$, $\psi_{M_2}^{(2)}$. For $M=|M_1|=|M_2|\ne0$, form the symmetrized and antisymmetrized products
\[
 \psi^+=\psi_M^{(1)}\psi_{-M}^{(2)}+\psi_{-M}^{(1)}\psi_M^{(2)},\qquad
 \psi^-=\psi_M^{(1)}\psi_{-M}^{(2)}-\psi_{-M}^{(1)}\psi_M^{(2)}.
\]
Reflection in a vertical plane (one containing the molecular axis) reverses the sign of the projection of angular momentum on the axis, so that $\psi_M^{(1)}$, $\psi_M^{(2)}$ pass respectively into $\psi_{-M}^{(1)}$, $\psi_{-M}^{(2)}$, and conversely.
Under this reflection, $\psi^+$ remains unchanged, whereas $\psi^-$ changes sign; hence the former corresponds to a $\Sigma^+$ term and the latter to a $\Sigma^-$ term.
Thus every value of $M$ gives one $\Sigma^+$ term and one $\Sigma^-$ term.
Since $M$ can take $L_2$ distinct values ($M=1,\ldots,L_2$), there are altogether $L_2$ terms of each type, $\Sigma^+$ and $\Sigma^-$.


If instead $M_1=M_2=0$, the molecular wavefunction has the form $\psi=\psi_0^{(1)}\psi_0^{(2)}$. To determine how $\psi_0^{(1)}$ behaves under reflection in a vertical plane, choose coordinates centered on the first atom, with the $z$ axis along the molecular axis. Reflection in the vertical plane $xz$ is equivalent to inversion through the coordinate origin followed by a rotation through $180^\circ$ about the $y$ axis. Inversion multiplies $\psi_0^{(1)}$ by $P_1$, where $P_1=\pm1$ is the parity of the given state of the first atom. The action on a wavefunction of an infinitesimal rotation, and hence of any finite rotation, is completely determined by the atom's total orbital angular momentum. It is therefore enough to consider the special case of a one-electron atom whose electron has orbital angular momentum $l$ (and angular-momentum component $m=0$ along $z$); replacing $l$ by $L$ in the result gives the desired answer for an arbitrary atom. Apart from a constant factor, the angular part of the electron wavefunction for $m=0$ is $P_l(\cos\theta)$ (see (28.8)). A rotation through $180^\circ$ about the $y$ axis gives $x\to-x$, $y\to y$, $z\to-z$, or, in spherical coordinates, $r\to r$, $\theta\to\pi-\theta$, $\varphi\to\pi-\varphi$. Consequently $\cos\theta\to-\cos\theta$, and $P_l(\cos\theta)$ is multiplied by $(-1)^l$.

\SourcePage{373}{376}
Thus reflection in a vertical plane multiplies $\psi_0^{(1)}$ by $(-1)^{L_1}P_1$. Similarly, $\psi_0^{(2)}$ is multiplied by $(-1)^{L_2}P_2$, so the net factor multiplying $\psi=\psi_0^{(1)}\psi_0^{(2)}$ is $(-1)^{L_1+L_2}P_1P_2$. The term is $\Sigma^+$ or $\Sigma^-$ according as this factor equals $+1$ or $-1$.

Collecting these results, of the total $(2L_2+1)$ $\Sigma$ terms (for each possible multiplicity), $L_2+1$ are $\Sigma^+$ terms and $L_2$ are $\Sigma^-$ terms if $(-1)^{L_1+L_2}P_1P_2=+1$; the numbers are interchanged if $(-1)^{L_1+L_2}P_1P_2=-1$.

We turn now to a molecule composed of identical atoms. The rules for combining the atomic spins and orbital angular momenta into the molecular values $S$ and $\Lambda$ remain the same as for a molecule made from different atoms. The new feature is that the terms may be even or odd. Two cases must be distinguished: the atoms may be in the same state or in different states.

If the atoms are in different states,\footnote{This includes, in particular, the combination of a neutral atom with an ionized one.} the total number of possible terms is doubled relative to the number for two different atoms. Indeed, inversion through the origin (placed at the midpoint of the molecular axis) interchanges the states of the two atoms. Symmetrizing or antisymmetrizing the molecular wavefunction with respect to this interchange gives two terms with the same $\Lambda$ and $S$, one even and the other odd. Thus the total numbers of even and odd terms are equal.

If, however, the two atoms are in identical states, the total number of states remains the same as for a molecule with dissimilar atoms. With regard to the parity of these states, an analysis (not reproduced here because it is cumbersome)\footnote{See E.~Wigner, E.~Witmer~// Zs. Physik. 1928. V.~51 S.~850.} gives the following results.


\SourcePage{374}{377}
Let $N_g$, $N_u$ denote the numbers of even and odd terms having specified values of $\Lambda$ and $S$. Then
\[
\begin{array}{ll}
\text{if $\Lambda$ is odd,}&N_g=N_u,\\
\text{if $\Lambda$ is even and $S$ is even $(S=0,2,4,\ldots)$,}&N_g=N_u+1,\\
\text{if $\Lambda$ is even and $S$ is odd $(S=1,3,\ldots)$,}&N_u=N_g+1.
\end{array}
\]
Finally, the $\Sigma$ terms must further be divided into $\Sigma^+$ and $\Sigma^-$. The rule here is
\[
\begin{array}{ll}
\text{if $S$ is even,}&N_g^+=N_u^-+1=L+1,\\
\text{if $S$ is odd,}&N_u^+=N_g^-+1=L+1
\end{array}
\]
(where $L_1=L_2\equiv L$). Every $\Sigma^+$ term has parity $(-1)^S$, and every $\Sigma^-$ term has parity $(-1)^{S+1}$.

Besides the relation discussed above between molecular terms and the terms of the atoms obtained as $r\to\infty$, one may also ask how the molecular terms are related to the terms of the ``united atom'' that would result as $r\to0$, when the two nuclei are brought to one point (for example, the relation between the terms of the $\mathrm H_2$ molecule and the He atom). The following rules are obtained immediately. A term of the ``united'' atom with spin $S$, orbital angular momentum $L$, and parity $P$ gives, when its constituent atoms are separated, molecular terms with spin $S$ and axial angular-momentum projection $\Lambda=0,1,\ldots,L$, one term for each of these values of $\Lambda$. The parity of a molecular term is the same as the parity $P$ of the atomic term ($g$ for $P=+1$ and $u$ for $P=-1$). The molecular term with $\Lambda=0$ is a $\Sigma^+$ term if $(-1)^LP=+1$, and a $\Sigma^-$ term if $(-1)^LP=-1$.

\ProblemHeading 1. Determine the possible terms of the molecules $\mathrm H_2$, $\mathrm N_2$, $\mathrm O_2$, and $\mathrm{Cl}_2$ that can result when atoms in their ground states combine.

\SolutionHeading The rules stated in the text give the following possible terms:
\[
\begin{array}{c@{\quad}l@{\quad}l}
\text{molecule }\mathrm H_2& (\text{atoms in }{}^2S\text{ states}):&{}^1\Sigma_g^+,\ {}^3\Sigma_u^+;\\
\mathrm N_2&(\text{atoms in }{}^4S\text{ states}):&{}^1\Sigma_g^+,\ {}^3\Sigma_u^+,\ {}^5\Sigma_g^+,\ {}^7\Sigma_u^+;\\
\mathrm{Cl}_2&(\text{atoms in }{}^2P\text{ states}):&2{}^1\Sigma_g^+,\ {}^1\Sigma_u^-,\ {}^1\Pi_g,\ {}^1\Pi_u,\ {}^1\Delta_g,\\
&&2{}^3\Sigma_u^+,\ {}^3\Sigma_g^-,\ {}^3\Pi_g,\ {}^3\Pi_u,\ {}^3\Delta_u;\\
\mathrm O_2&(\text{atoms in }{}^3P\text{ states}):&2{}^1\Sigma_g^+,\ {}^1\Sigma_u^-,\ {}^1\Pi_g,\ {}^1\Pi_u,\ {}^1\Delta_g,\\
&&2{}^3\Sigma_u^+,\ {}^3\Sigma_g^-,\ {}^3\Pi_u,\ {}^3\Pi_g,\ {}^3\Delta_u,\\
&&2{}^5\Sigma_g^+,\ {}^5\Sigma_u^-,\ {}^5\Pi_g,\ {}^5\Pi_u,\ {}^5\Delta_g.
\end{array}
\]
(A numeral before a term symbol gives the number of terms of that type when the number exceeds one.)

\ProblemHeading 2. Do the same for the molecules HCl and CO.

\SourcePage{375}{378}
\SolutionHeading For unlike atoms, the parities of their states must also be taken into account. Formula (31.6) shows that the ground states of H, O, and C are even, while that of Cl is odd (for the electronic configurations of the atoms, see Table~3). The rules given in the text yield
\[
\begin{array}{c@{\quad}l@{\quad}l}
\text{molecule HCl}&(\text{atoms in }{}^2S_g,{}^2P_u\text{ states}):&{}^{1,3}\Sigma^+,\ {}^{1,3}\Pi;\\
\text{molecule CO}&(\text{atoms in }{}^3P_g,{}^3P_g\text{ states}):&2{}^{1,3,5}\Sigma^+,\ {}^{1,3,5}\Sigma^-,\\
&&2{}^{1,3,5}\Pi,\ {}^{1,3,5}\Delta.
\end{array}
\]
